\section{Discussion}

The contractual runtime trades additional monitoring and trace overhead for improved diagnosability and bounded recovery. The key limitation is that semantic checks require domain-specific verifiers or benchmark-specific policies. The harness is designed so stronger verifiers can be added without changing workload data.

\subsection{Cost and Reliability}

Contracts and traces add runtime work. The relevant tradeoff is not raw latency alone, but cost per successful task under fault injection. A retry-only strategy may appear simple, but can amplify tool and LLM calls without identifying the contaminated artifact. \system pays monitoring overhead to make violations and recovery decisions inspectable. In the current main matrix, schema-only slightly leads raw task success, while the full runtime is the only configuration with non-zero bounded recovery and lower cost per successful task than retry-only. This is a failure-containment claim, not a task-success leaderboard claim.

\subsection{Limits of the Current Artifact}

The current artifact is external-benchmark backed: all four runtime methods are driven by the \texttt{SyncLLMClient} against public benchmark task fixtures, and the main matrix is gated by an audit script that requires both external main-matrix rows and AgentChangeBench rows. The harness still includes deterministic smoke modes for development and regression tests, but those rows are not used for the paper's main claims. The remaining limitations are important: the scorer is fixture-backed rather than a full benchmark-native simulator rerun; the matrix uses a single Qwen3-8B endpoint; and several trace metrics such as mean time to detect, mean time to recover, fault propagation depth, and contaminated artifact count are placeholders unless timestamped traces and artifact-propagation annotations are present. These limits motivate a controlled successful-trajectory fault-injection study as the next experiment.

\subsection{Threats to Validity}

Fault distributions may not match all production incidents. Semantic contracts can be incomplete or domain-biased. The task fixtures are public and external, but the execution uses one live LLM endpoint and a constrained sampling plan; stronger production claims require benchmark-native simulator reruns, additional models, larger task sweeps, independent reruns, and controlled fault injection on trajectories that first succeed without injected faults. These limits affect generality, not the artifact's audit trail: the JSONL traces, summary tables, figures, and audit report are all generated from the same external main-matrix directory.
